\exercisesfor{3}

The letter X denotes a set.

\begin{enumerate}
\def\labelenumi{\arabic{enumi}.}
\tightlist
\item
  \begin{enumerate}
  \def\labelenumii{(\alph{enumii})}
  \tightlist
  \item
    Show that every element \(u \in \mathrm{A}^{+}(\mathrm{X})\) can be written uniquely in the form \(u = \sum_{x \in \mathrm{X}} u_{x} x\) , where \(u_{x} \in \mathrm{A}(\mathrm{X})\) .
  \end{enumerate}
\end{enumerate}

\begin{enumerate}
\def\labelenumi{(\alph{enumi})}
\setcounter{enumi}{1}
\tightlist
\item
  Show that \(\{0\}\) is the only submodule of \(\mathrm{A}^{+}(\mathrm{X})\) which is stable under all the mappings \(u \mapsto u_{x} (x \in \mathrm{X})\) . (If a is such a submodule and \(a \neq \{0\}\) , consider a non-zero element of a of minimal degree.)
\end{enumerate}

\begin{enumerate}
\def\labelenumi{\arabic{enumi}.}
\setcounter{enumi}{1}
\tightlist
\item
  Let g be a Lie algebra which is a free module; it is identified by means of \(\sigma: g \to Ug\) with its image in the enveloping algebra \(Ug\) . Let \(U^{+}g\) be the kernel of the canonical homomorphism \(Ug \to K\) . Let i be a mapping of X into g such that \(i(X)\) generates g as a Lie algebra.
\end{enumerate}

\begin{enumerate}
\def\labelenumi{(\alph{enumi})}
\item
  Show that \(\mathbf{U}^{+}\mathfrak{g}\) is generated by \(i(\mathbf{X})\) as a left Ug-module.
\item
  Show the equivalence of the following properties:
\item
  g is free with basic family i: X → g.
\end{enumerate}

\begin{enumerate}
\def\labelenumi{(\roman{enumi})}
\setcounter{enumi}{1}
\tightlist
\item
  The family \(i\) is a basis of the left Ug-module \(\mathbf{U}^{+}\mathfrak{g}\) .
\end{enumerate}

(The implication (i) \(\Rightarrow\) (ii) follows from Exercise 1 (a) using the isomorphism \(\mathrm{Ug} \to \mathrm{A}(\mathrm{X})\) . To show (ii) \(\Rightarrow\) (i), apply Exercise 1 (b) to the kernel of the homomorphism \(\mathrm{A}(\mathrm{X}) \to \mathrm{Ug}\) defined by \(i\) .)

\begin{enumerate}
\def\labelenumi{\arabic{enumi}.}
\setcounter{enumi}{2}
\tightlist
\item
  Let \(x \in \mathbf{X}\) . Show that the centralizer of \(x\) in \(\mathrm{A}(\mathbf{X})\) is the subalgebra generated by \(x\) . Deduce that the only elements of \(\mathrm{L}(\mathbf{X})\) which commute with \(x\) are the multiples of \(x\) . In particular, the centre of \(\mathrm{L}(\mathbf{X})\) is (0) if \(\operatorname{Card}(\mathbf{X}) \geqslant 2\) and the centre of \(\mathrm{L}(\mathbf{X}) / \left( \sum_{n > p} \mathrm{L}^n(\mathbf{X}) \right)\) reduces to the canonical image of \(\mathrm{L}^p(\mathbf{X})\) .
\end{enumerate}


¶ 4. Suppose that K is a field of characteristic p \textgreater{} 0.

\begin{enumerate}
\def\labelenumi{(\alph{enumi})}
\item
  Let \(\mathbf{H}\) be a basis of the Lie subalgebra \(\mathbf{L}(\mathbf{X})\) of \(\mathbf{A}(\mathbf{X})\) . Show (using the isomorphism \(\mathbf{U}(\mathbf{L}(\mathbf{X})) \to \mathbf{A}(\mathbf{X})\) and the Poincaré-Birkhoff-Witt Theorem) that the elements \(h^{p^n}\) ( \(h \in \mathbf{H}\) , \(n \in \mathbf{N}\) ) are linearly independent over \(\mathbf{K}\) .
\item
  If \(m\) is an integer \(\geqslant 0\) , let \(L_m\) denote the submodule of \(A(X)\) with basis the \(h^{p^n}\) , where \(h \in H, n \leqslant m\) . Show that \(L_m\) is a Lie subalgebra of \(A(X)\) and that, if \(a \in L_m - 1\) , then \(a^p \in L_m\) (argue by induction on \(m\) and use the Jacobson formulae, cf.~Chapter I, § 1, Exercise 19). Deduce that \(L_m\) does not depend on the choice of \(H\) and that it is the Lie subalgebra of \(A(X)\) generated by the \(x^{p^n}\) , where \(x \in X, n \leqslant m\) .
\item
  Let \(\mathbf{L}(\mathbf{X}, p)\) be the union of the \(\mathbf{L}_m\) for \(m \geqslant 0\) . Show that \(\mathbf{L}(\mathbf{X}, p)\) is the smallest Lie \(p\) -subalgebra of \(\mathbf{A}(\mathbf{X})\) containing \(\mathbf{X}\) (cf.~Chapter I, §1, Exercise 20); it admits as basis the family of \(h^{p^n}\) , where \(h \in \mathbf{H}\) , \(n \in \mathbf{N}\) .
\item
  Let \(\tilde{\mathbf{U}}(\mathbf{L}(\mathbf{X},p))\) be the restricted enveloping algebra of \(\mathbf{L}(\mathbf{X},p)\) , cf.~Chapter I, § 2, Exercise 6. Show that the injection of \(\mathbf{L}(\mathbf{X},p)\) into \(\mathbf{A}(\mathbf{X})\) can be extended to an isomorphism of \(\tilde{\mathbf{U}}(\mathbf{L}(\mathbf{X},p))\) onto \(\mathbf{A}(\mathbf{X})\) . (Use the Poincaré-Birkhoff-Witt Theorem and the exercise quoted above.) Deduce that \(\mathbf{L}(\mathbf{X},p)\) is the set of primitive elements of \(\mathbf{A}(\mathbf{X})\) , cf.~§ 1, Exercise 12.
\item
  Let \(f\) be a mapping of \(\mathbf{X}\) into a Lie \(p\) -algebra \(\mathfrak{g}\) . Show that \(f\) can be extended uniquely to a \(p\) -homomorphism \(\mathbf{F} \colon \mathbf{L}(\mathbf{X}, p) \to \mathfrak{g}\) . (Begin by extending \(f\) to an algebra homomorphism of \(\mathrm{A}(\mathbf{X})\) into \(\tilde{\mathrm{U}}\mathfrak{g}\) .)
\end{enumerate}

(The Lie \(p\) -algebra \(\mathbf{L}(\mathbf{X}, p)\) is called the free Lie \(p\) -algebra on the set \(\mathbf{X}\) .)

